\chapter{Réalisation de Palatini–Cartan–Holst du paramètre de Barbero–Immirzi et transduction physique du spectre d’aire}
\label{chap:parte-iv-68}

La dérivation structurelle du paramètre et du spectre d’aire précède ce chapitre. On construit ici la tétrade, la connexion, la courbure et l’action du premier ordre, et ce n’est qu’ensuite que le paramètre de Holst est spécialisé comme réalisation physique du lecteur aréal.

\section{Courbure, torsion et action d’Einstein--Cartan--Holst}
\label{chap:gravedad-cartan-holst}

La structure de transport HMT distingue la variable de mémoire, l’holonomie et le défaut de fermeture. Sa réalisation gravitationnelle se formule au moyen d’une tétrade et d’une connexion. Cette section établit la géométrie du premier ordre et spécifie le type de l’application entre celle-ci et l’espace des états TPK\@.

\subsection{Identité discrète de Gauss–Stokes et transport de l’incidence vers la frontière}

Soit \(K\) un complexe cellulaire, soit \(W\) un groupe abélien et soit \(\Omega\in C^1(K;W)\) une cochaîne. Nous notons \(\delta:C^1(K;W)\to C^2(K;W)\) le cobord cellulaire et \(\langle\,\cdot\,,\,\cdot\,\rangle\) l’accouplement entre cochaînes à valeurs dans \(W\) et chaînes cellulaires entières.

\begin{theorem}[Identité cellulaire universelle de Stokes]
Pour toute \(2\)-chaîne finie \(z\in C_2(K;\mathbb Z)\),
\[
 \boxed{\langle\delta\Omega,z\rangle
 =\langle\Omega,\partial z\rangle.}
\]
\end{theorem}
\begin{proof}
C’est l’identité qui définit le cobord comme opérateur adjoint du bord cellulaire. La finitude de \(z\) garantit que les deux accouplements sont des sommes finies dans \(W\).
\end{proof}

\begin{corollary}[Gauss--Stokes sur une pseudovariété cellulaire]
Soit \(K\) une pseudovariété cellulaire régulière, bidimensionnelle, finie et compacte, orientée de manière cohérente, éventuellement à bord. Si \(z_K\) est la somme de ses \(2\)-cellules avec l’orientation cohérente, alors
\[
 \sum_{f\in K^{(2)}}(\delta\Omega)(f)
 =\sum_e[\partial z_K:e]\,\Omega(e).
\]
En particulier, la régularité permet d’identifier les incidences non annulées de \(\partial z_K\) aux arêtes du bord géométrique.
\end{corollary}
\begin{proof}
On applique le théorème à \(z_K\). L’orientation cohérente fait apparaître chaque arête intérieure deux fois dans \(\partial z_K\), avec des coefficients opposés ; les incidences non annulées forment la chaîne de bord.
\end{proof}

La version non abélienne ne s’obtient pas en remplaçant simplement les sommes par des produits. Dans cet ouvrage, elle n’est utilisée que lorsque les transports vers une base commune, un ordre des holonomies et la conjugaison correspondante ont été fixés. En dehors de ce sous-domaine, aucune identité de Stokes non abélienne n’est affirmée.

\subsection{Capacité aréale de la cellule qutrit : \(81\log 3\) et transport collectif d’incidence}

\begin{theorem}[Coupe minimale d’un cube discret]
Dans le graphe cubique \(N\times N\times N\) de capacité unitaire, toute coupe séparant deux faces opposées a une capacité au moins égale à \(N^2\), et il en existe une de capacité \(N^2\).
\end{theorem}
\begin{proof}
Il existe \(N^2\) chemins disjoints en arêtes qui traversent le cube dans la direction normale aux deux faces. Toute coupe doit les intercepter, donc sa capacité est au moins \(N^2\). Une couche transversale contient exactement \(N^2\) arêtes et atteint la borne.
\end{proof}

Pour \(N=9\), la coupe vaut 81 et
\[
 9^3=729=27^2.
\]
La dernière égalité permet de regrouper six trits sous forme d’un volume \(9^3\) ou d’une grille de cardinalité \(27^2\). C’est une bijection d’ensembles ; ce n’est pas un isomorphisme entre \((\ZZ/9\ZZ)^3\) et \((\ZZ/27\ZZ)^2\), dont les exposants diffèrent.

\subsection{Tétrade, connexion et \(2\) défauts}

Soient \(e^I\) une tétrade et \(\omega^{IJ}=-\omega^{JI}\) une connexion de Lorentz. Définissons
\[
 T_{\rm tor}^I=D_\omega e^I
 =\dd e^I+\omega^I{}_J\wedge e^J,
\]
\[
 F^{IJ}=\dd\omega^{IJ}+\omega^I{}_K\wedge\omega^{KJ},
 \qquad
 g=\eta_{IJ}e^I\otimes e^J.
\]
La décomposition
\[
 \omega^{IJ}=\mathring\omega^{IJ}(e)+K^{IJ}
\]
sépare la connexion de Levi--Civita et la contorsion. Alors
\[
 T_{\rm tor}^I=K^I{}_J\wedge e^J.
\]
Courbure et torsion sont des défauts différents : \(F\) mesure la fermeture du transport des vecteurs ; \(T_{\rm tor}\), la fermeture de la tétrade. La mémoire \(T\) du principe d’Ambivalence n’est pas non plus une torsion géométrique : seul un foncteur de réalisation peut les relier.

\subsection{Action du 1er ordre}

Soit \(M\) une variété lisse orientée de dimension quatre. Si \(\partial M\ne\varnothing\), nous supposons des variations à support compact dans l’intérieur ou, de manière équivalente pour le problème variationnel considéré, un terme de bord et des conditions aux limites annulant la contribution frontalière. Lorsque \(\Psi\) contient des champs spinoriels, nous supposons en outre une structure de spin et des champs compatibles avec elle. Soient
\[
 \Sigma^{IJ}=e^I\wedge e^J,
 \qquad
 (\star X)^{IJ}=\frac12\epsilon^{IJ}{}_{KL}X^{KL},
 \qquad
 \kappa=\frac{8\pi G}{c^4},
 \qquad
 \gamma\in\mathbb R\setminus\{0\}.
\]
L’action de Palatini--Cartan--Holst est
\[
\boxed{
 S[e,\omega,\Psi]
 =\frac1{2\kappa c}\int_M
 \Sigma_{IJ}\wedge\left(\star F^{IJ}+\frac1\gamma F^{IJ}\right)
 -\frac\Lambda{\kappa c}\int_M\operatorname{vol}_e
 +S_{\rm m}[e,\omega,\Psi].}
\]
La variation de cette action détermine les corrections torsionnelles \cite{Sciama1964,Kibble1961,Hehl1976,Holst1996}.

Définissons le courant de spin par
\[
 \delta_\omega S_{\rm m}
 =-\frac12\int_M\delta\omega^{IJ}\wedge\sigma_{IJ}.
\]

\begin{theorem}[Équation de Cartan--Holst]
Sous les hypothèses géométriques et de bord précédentes, la variation par rapport à \(\omega\) donne
\[
 D_\omega\!\left[(\star+\gamma^{-1}I)\Sigma^{IJ}\right]
 =\kappa c\,\sigma^{IJ},
\]
avec la normalisation globale fixée par la définition précédente de \(\sigma^{IJ}\).
\end{theorem}
\begin{proof}
On utilise \(\delta F^{IJ}=D_\omega\delta\omega^{IJ}\), on intègre par parties et l’on annule le coefficient de la variation arbitraire \(\delta\omega^{IJ}\). Avec la convention de signe fixée dans la définition de \(\sigma^{IJ}\), le terme de matière passe au membre de droite avec un signe positif. Le facteur \(1/(2\kappa c)\) produit le coefficient \(\kappa c\). Le terme cosmologique ne dépend pas de \(\omega\), et les conditions aux limites éliminent le terme engendré par l’intégration par parties.
\end{proof}

En signature lorentzienne, \(\star^2=-I\) sur les bivecteurs et, pour \(\gamma\in\RR\setminus\{0\}\),
\[
 (\star+\gamma^{-1}I)^{-1}
 =\frac{\gamma^2}{1+\gamma^2}(-\star+\gamma^{-1}I).
\]

\begin{corollary}[Découplage torsionnel]
Pour une cotétrade non dégénérée et \(\gamma\in\mathbb R\setminus\{0\}\),
\[
 \sigma^{IJ}=0\Longrightarrow T_{\rm tor}^I=0
 \Longrightarrow\omega=\mathring\omega(e).
\]
\end{corollary}

Pour expliciter la normalisation, définissons les sources normalisées par l’action au moyen de
\[
 \delta_g S_{\rm m}
 =-\frac12\int_M\operatorname{vol}_e\,
 \mathcal T^{S}_{\mu\nu}\,\delta g^{\mu\nu}.
\]
Lorsque les coordonnées de la tétrade ont la dimension d’une longueur, le tenseur énergie--impulsion physique conventionnel est \(T^{\rm phys}_{\mu\nu}:=c\,\mathcal T^{S}_{\mu\nu}\). Notons de même \(\mathcal U^{S,\rm spin}\) la contribution normalisée par l’action et \(U^{\rm phys,\rm spin}:=c\,\mathcal U^{S,\rm spin}\) sa convention physique. En éliminant algébriquement la connexion, l’équation métrique prend alors les deux formes équivalentes
\[
 G_{\mu\nu}+\Lambda g_{\mu\nu}
 =\kappa c\,
 \bigl(\mathcal T_{\mu\nu}^{S,\rm LC}
       +\mathcal U_{\mu\nu}^{S,\rm spin}\bigr)
 =\kappa\,
 \bigl(T_{\mu\nu}^{\rm phys,LC}
       +U_{\mu\nu}^{\rm phys,spin}\bigr),
\]
où la contribution du spin est quadratique dans le courant correspondant et dépend du couplage de matière. Le coefficient de cette contribution ne peut être fixé sans spécifier l’action fermionique et les conventions de dualité.

\subsection{Identité de Bianchi avec torsion et bilan de flux dans la connexion de Cartan}

Pour toute connexion,
\[
 D_\omega T_{\rm tor}^I=F^I{}_J\wedge e^J,
 \qquad
 D_\omega F^{IJ}=0.
\]
Après élimination de la torsion,
\[
\nabla^\mu(T_{\mu\nu}^{\rm phys,LC}
          +U_{\mu\nu}^{\rm phys,spin})=0.
\]
Ces identités sont des tests de cohérence obligatoires : tout terme gravitationnel HMT doit provenir d’une action compatible ou satisfaire le bilan de Noether correspondant.

\subsection{Application de réalisation et spécialisation du paramètre de Holst}

L’insertion définie dans \cref{mdg:chap:barbero} prend
\[
 \gamma=\Gamma_{6\to8}.
\]
Une fois cette substitution effectuée, toutes les conséquences variationnelles précédentes sont des théorèmes de l’action spécialisée. La sélection physique de cette action par le transport HMT s’exprime au moyen d’une application de réalisation
\[
 \mathfrak R_{\rm grav}:\mathcal X_{\rm TPK}
 \dashrightarrow\{(e,\omega,\Psi)\}
\]
dotée de trois propriétés : entrelacer le transport TPK avec celui de \(\omega\), identifier les classes d’holonomie et fournir les échelles de \(G\) et de \(\Lambda\). Ces propriétés fixent le domaine précis d’une réalisation gravitationnelle du formalisme.

Pour une cotétrade non dégénérée, \(\gamma\) réel non nul et une matière sans courant de spin, le secteur torsionnel se réduit à la formulation du premier ordre de la relativité générale avec la même constante cosmologique, sous les conditions aux limites déjà déclarées. Dans un fluide de spin, la loi d’échelle \(s^2\propto a^{-6}\) peut produire un rebond d’Einstein--Cartan ; obtenir sa densité et ses conditions initiales à partir de l’état TPK est une question distincte de la cohérence variationnelle de l’action.


\section{Retour nonadique, groupoïde de mémoire et condition d’entrelacement de Cartan}

La hiérarchie \(27\to54\to108\), le recouvrement cumulatif dirigé et le groupoïde de mémoire fournissent la condition exacte sous laquelle le retour du TPK admet une connexion de Cartan compatible. Ce pont précède la spécialisation de Palatini--Cartan--Holst et conserve intégralement son domaine et ses critères de réfutation.

\begingroup
\let\chapter\section
\let\section\subsection
\let\subsection\subsubsection
\let\subsubsection\paragraph
\let\clearpage\relax
\let\cleardoublepage\relax
\chapter{Retour nonadique, mémoire et holonomie de Cartan}
\label{chap:puente-retorno-holonomia-cartan}
\index{retour!27--54--108}\index{holonomie!TPK--Cartan}
\index{mémoire!distinction par rapport à la torsion}

La réalisation gravitationnelle ne commence pas par identifier une variable discrète
à une grandeur géométrique. Elle commence par distinguer trois niveaux : le retour
du transport observable, la mémoire qui différencie des histoires ayant la même
projection et l'holonomie d'une connexion. Le pont formulé ici
reçoit les retours et les cocycles déjà démontrés dans HMT ; il ne les reconstruit pas
à partir de l'action de Holst, d'une masse ou d'une constante métrologique.

\section{Hiérarchie exacte des retours}
\label{sec:jerarquia-retornos-27-54-108}

Soit \(F_{\rm obs}\) la chronologie observable définie dans le noyau HMT et soit
\[
 \mathcal K_{27}:=F_{\rm obs}^{27}.
\]
Le symbole \(\mathcal K_{27}\) désigne exclusivement la composition de
vingt-sept mises à jour de la boucle positionnelle et orientée ; il ne désigne pas une
fenêtre générique, une trace de \(108\) pas ni l'état TPK enrichi
complet. Notons \(p\) la projection sur les positions des deux
curseurs et \(\mathsf J_{\rm or}\) l'inversion simultanée de leurs
orientations. La démonstration primaire se trouve dans
\cref{prop:retorno-fobs-identidad} ; ses conséquences typées sont
\[
 p(\mathcal K_{27}x)=p(x),\qquad
 o(\mathcal K_{27}x)=\mathsf J_{\rm or}\,o(x),
 \qquad \mathsf J_{\rm or}^{\,2}=I,
\]
\[
 F_{\rm obs}^{54}:
 \text{rétablit les positions et les orientations},
 \qquad
 F_{\rm obs}^{108}=I:
 \text{rétablit aussi la phase observable}.
\]
Par conséquent, retour de position, retour d'orientation et retour de phase sont
des affirmations différentes. Aucune d'elles n'implique à elle seule le retour de
tous les champs d'une extension enrichie.

Le quotient de modes et les cocycles d'orientation sont établis dans
\cref{sec:cociclos-modo-orientacion,thm:descenso-necesario-fav}. Dans l'ordre
surfacique--volumétrique,
\[
 N_{27}=9F_{\rm av},
 \qquad
 F_{\rm av}=
 \begin{pmatrix}0&1\\1&1\end{pmatrix}.
\]
Si
\[
 \Omega_{\rm sw}(\gamma)=\sum_{a\subset\gamma}\omega_{\rm sw}(a),
 \qquad
 S(\gamma)=\sum_{a\subset\gamma}|\omega_{\rm sw}(a)|,
 \qquad
 G_{\rm or}(\gamma)=\sum_{a\subset\gamma}g(a),
\]
alors un bloc de vingt-sept pas contient trois tours nonadiques et une
inversion d'orientation :
\[
 \Omega_{\rm sw}(\mathcal K_{27})=0,\qquad
 S(\mathcal K_{27})=18,\qquad
 G_{\rm or}(\mathcal K_{27})=1.
\]
Dans le retour observable de \(108\) pas,
\[
 \Omega_{\rm sw}(C_{108})=0,\qquad
 S(C_{108})=72,\qquad
 G_{\rm or}(C_{108})=4.
\]
Ainsi, pour la fonctionnelle additive déclarée
\[
 \mathcal A_\ast(\gamma)
 :=G_{\rm or}(\gamma)+\lambda_\ast S(\gamma),
\]
on obtient, sans introduire d'échelle dimensionnelle,
\[
 \mathcal A_\ast(\mathcal K_{27})=1+18\lambda_\ast,
 \qquad
 \mathcal A_\ast(C_{108})=4+72\lambda_\ast.
\]
Ces valeurs sont des comptages du registre. Leur utilisation thermique ou
gravitationnelle ultérieure n'altère pas leur provenance.

\subsection{Système dirigé de revêtements, raffinement et limite de la connexion de Cartan}
\label{sec:cobertura-acumulativa-dirigida}

Soit \(\mathcal O_{108}\) l'orbite observable et, pour chaque arête chronologique
\(e\), soit
\[
 c(e)=\bigl(g(e),s(e)\bigr)\in\mathbb N^2
\]
le couple formé par les indicateurs d'inversion d'orientation et de changement
effectif de feuillet. Pour une histoire dirigée
\(\gamma=e_1\cdots e_n\), on définit
\[
 c(\gamma)=\sum_{j=1}^{n}c(e_j).
\]
La loi de concaténation prend la forme
\[
 c_{m+n}(x)
 =c_m(x)+c_n\!\left(F_{\rm obs}^{m}x\right).
\]
Les comptages précédents expriment, uniformément sur l'orbite,
\[
 c_{27}(x)=(1,18),
 \qquad
 c_{108}(x)=(4,72).
\]

\begin{theorem}[Revêtement cumulatif de la chronologie observable]
\label{thm:cobertura-acumulativa-dirigida}
Sur
\[
 \widetilde{\mathcal O}_{108}^{+}
 :=\mathcal O_{108}\times\mathbb N^2
\]
définissons
\[
 \widetilde F^{+}(x,m)
 :=\bigl(F_{\rm obs}x,m+c(x\to F_{\rm obs}x)\bigr).
\]
Alors le transport de vingt-sept pas satisfait
\[
 \widetilde{\mathcal K}_{27}^{+}(x,m)
 =\bigl(\mathcal K_{27}x,m+(1,18)\bigr),
\]
et
\[
 \left(\widetilde{\mathcal K}_{27}^{+}\right)^4(x,m)
 =\bigl(x,m+(4,72)\bigr)\neq(x,m).
\]
Par conséquent, sa projection observable est d'ordre quatre, tandis que
l'action cumulative du monoïde chronologique possède un ordre de semi-groupe infini.
\end{theorem}

\begin{proof}
La loi cocyclique et la constance de \(c_{27}\) donnent
\[
 c_{108}(x)
 =\sum_{j=0}^{3}c_{27}(\mathcal K_{27}^{j}x)
 =4(1,18)=(4,72).
\]
La première coordonnée revient par
\(\mathcal K_{27}^{4}=I\) ; la seconde ne revient pas. Après \(q>0\) itérations,
l'incrément est \(q(4,72)\), non nul.
\end{proof}

\begin{remark}[Type de l'extension cumulative]
La construction précédente étend les histoires dirigées au moyen de compteurs
positifs. Ce n'est pas un relèvement automorphe du groupoïde
\(\mathcal G_{\rm mem}\), et elle ne démontre pas que la coordonnée illimitée appartienne
déjà à l'état \(\mathcal X_{\mathrm{TPK}}^{\mathrm{enr}}\). Sa complétion de Grothendieck
\(\mathbb Z^2\) peut être utilisée pour l'analyse harmonique, mais elle ne transforme pas
le coût positif
\(\mathcal A_\ast=G_{\rm or}+\lambda_\ast S\) en un caractère du groupe
d'isotropie : lorsqu'on inverse formellement une voie, ses indicateurs ne prennent pas
le signe opposé.
\end{remark}

\paragraph{Domaine de validité de l'extension cumulative.}
La hiérarchie \(27\to54\to108\), les incidences de \(F_{\rm av}\) et les
cocycles précédents sont des identités exactes dans la chronologie et les
lecteurs déclarés. L'affirmation est restreinte à l'état observable que
ces lecteurs conservent ; l'étendre à l'état enrichi complet exigerait
de prouver le retour de chaque champ oublié.

\section{Groupoïde de mémoire}
\label{sec:grupoide-memoria-tpk}

Soit \(\mathsf P_{\rm TPK}^{\rm enr}\) le groupoïde des chemins enrichis :
ses objets sont des états TPK survivants et ses flèches sont des voies admissibles
avec origine et destination déclarées, composition par concaténation et inversion
par renversement de voie. Soit \(\sim_{\rm mem}\) une équivalence qui
n'identifie deux flèches que si elle conserve les champs de mémoire requis par
le lecteur considéré et qui est congruente avec les identités, l'inversion et la
concaténation. Nous définissons
\[
 \mathcal G_{\rm mem}
 :=\mathsf P_{\rm TPK}^{\rm enr}/\!\sim_{\rm mem}.
\]
Cette définition ne transforme pas la mémoire en torsion. La mémoire appartient à
l'état discret et distingue des histoires ; la torsion appartient à une connexion
géométrique et mesure \(D_\omega e\). Seule une application entre les deux domaines peut
les relier.

Une réalisation géométrique exige, comme données supplémentaires, une longueur
\(\ell_0>0\), une représentation du groupoïde
\[
 \rho_{\rm C}:
 \mathcal G_{\rm mem}
 \longrightarrow
 \operatorname{Spin}(1,3)\ltimes\RR^{1,3},
\]
et une application de champs
\[
 \mathfrak R_{\rm grav}:
 \mathsf P_{\rm TPK}^{\rm enr}
 \dashrightarrow
 \mathsf{Cartan}_{1,3}.
\]
Pour définir la seconde application, il faut spécifier son action sur les objets et les
voies, sa compatibilité avec la concaténation et l'inversion et la régularité de la
connexion résultante. Aucune de ces propriétés ne se déduit du simple comptage
des retours.

\section{Connexion conjointe et condition d'entrelacement}
\label{sec:conexion-cartan-conjunta}

Sous les hypothèses géométriques du \cref{chap:gravedad} ---variété
orientée de dimension quatre, cotétrade \(e^I\), connexion de Lorentz
\(\omega^{IJ}\) et conditions au bord qui y sont déclarées---, la connexion de
Cartan conjointe s'écrit
\[
 \mathcal A_{\ell_0}
 =
 \begin{pmatrix}
  \omega&\ell_0^{-1}e\\
  0&0
 \end{pmatrix}.
\]
Sa courbure est l'identité algébrique
\[
 \mathcal F_{\ell_0}
 =\dd\mathcal A_{\ell_0}
  +\mathcal A_{\ell_0}\wedge\mathcal A_{\ell_0}
 =
 \begin{pmatrix}
  F_\omega&\ell_0^{-1}T_{\rm tor}\\
  0&0
 \end{pmatrix},
\]
où
\[
 F_\omega^{IJ}
 =\dd\omega^{IJ}+\omega^I{}_K\wedge\omega^{KJ},
 \qquad
 T_{\rm tor}^{I}=D_\omega e^I.
\]
La courbure et la torsion apparaissent comme des composantes distinctes d'un même défaut
de transport. L'égalité matricielle découle algébriquement des données
explicites \(e,\omega,\ell_0\), mais ne construit pas ces entrées à partir de TPK.

Le contenu fort du pont est la commutativité des holonomies :
\[
 \boxed{
 \Hol_{\mathcal A_{\ell_0}}
 \bigl(\mathfrak R_{\rm grav}(\gamma)\bigr)
 =
 \rho_{\rm C}\!\left(
 \Hol_{\rm TPK}(\gamma)
 \right)}
 \qquad
 (\gamma\in\mathsf P_{\rm TPK}^{\rm enr}).
\]
L'égalité doit être compatible avec les identités, l'inversion, la concaténation,
la subdivision, le changement de point de base et la conjugaison. Alors seulement la mémoire
discrète est représentée dans une connexion ; elle n'est pas rebaptisée
torsion ou courbure.

\begin{criterion}[Conditions d'existence du pont TPK--Cartan]
\label{crit:puente-tpk-cartan}
Il n'existe pas de réalisation possédant la commutativité et la compatibilité déclarées
si l'un quelconque des faits suivants se produit :
\begin{enumerate}
 \item deux représentants équivalents d'une même voie produisent
 des holonomies géométriques non conjuguées ;
 \item l'image ne préserve pas les identités, les inverses ou la concaténation ;
 \item une subdivision admissible altère l'holonomie ;
 \item l'égalité des holonomies échoue pour \(\mathcal K_{27}\) ou pour ses
 concaténations de profondeurs \(54\) et \(108\) ;
 \item l'application identifie mémoire et torsion sans exhiber un morphisme qui
 préserve leurs types ;
 \item un choix descendant ---l'action de Holst, \(G\), une masse ou
 une carte SI--- est utilisé pour redéfinir \(A\), \(C^\ast\) ou le retour
 TPK d'origine.
\end{enumerate}
\end{criterion}

\paragraph{Antériorité de la fonctionnelle \(\Gamma_{6\to8}\) par rapport à la réalisation Palatini--Cartan--Holst.}
Le \cref{mdg:chap:barbero} fixe, une fois \(A,C^\ast,q_\pm\) construits, la
fonctionnelle interne \(\Gamma_{6\to8}\). Le \cref{chap:gravedad} étudie ensuite
l'action Palatini--Cartan--Holst et ses équations variationnelles. Le présent pont
n'introduit pas de voie inverse entre les deux étapes.

\endgroup

% Owner íntegro de la realización discreta del retorno, la holonomía y la
% respuesta colectiva. Se sitúa después del puente que declara sus hipótesis.
\begin{HMTScientificOwnerSubordinated}
\chapter{Gravité : holonomie, torsion et réponse collective}
\label{chap:gravedad}
\index{gravité}\index{Cartan--Holst}\index{holonomie}
\index{torsion}\index{graviton!non primitif}

La gravité n'est pas introduite ici au moyen d'une formule décimale pour \(G\).
Sa construction réunit deux branches déjà typées. Dans la première, une voie conserve
orientation et mémoire ; une connexion compare des repères en des points distincts ; son
holonomie mesure le retour et sa courbure, le défaut de fermeture. Dans la seconde,
la carte dimensionnelle construit \(G_{\rm int}(\theta_G)\) et
\(\kappa_{\rm int}\) à partir de l'échelle interne d'action, de la vitesse et
de la durée élémentaire. Ce n'est qu'après avoir disposé des deux branches que l'on écrit
l'action Palatini--Cartan--Holst. Aucune valeur métrologique externe ne participe à
cette antériorité.

\section{Retours : position, orientation, phase et mémoire}

Soit \(F_{\rm obs}\) la chronologie observable du transport et
\[
 \mathcal K_{27}=F_{\rm obs}^{27}.
\]
La hiérarchie exacte distingue trois retours :
\[
 \begin{array}{c|c}
 27\ \text{pas}&
 \text{retour de position et inversion d’orientation},\\
 54\ \text{pas}&
 \text{retour de position et d’orientation},\\
 108\ \text{pas}&
 \text{retour complet de la phase observable}.
 \end{array}
\]
Le retour observable n'efface pas le registre accumulé. Si \(g(e)\) compte
une inversion d'orientation et \(s(e)\) un changement effectif de feuillet, alors
\[
 c(\gamma)=\sum_{e\subset\gamma}(g(e),s(e))
\]
satisfait la loi cocyclique de concaténation. Dans les blocs canoniques,
\[
 c_{27}=(1,18),\qquad c_{108}=(4,72).
\]
Par conséquent, la projection observable de \(\mathcal K_{27}\) est d'ordre
quatre, tandis que la coordonnée cumulative continue de croître. Revenir à la
même phase ne signifie pas revenir à la même histoire.

\begin{theorem}[Revêtement cumulatif du retour]
\label{thm:cobertura-retorno-grav-rev6}
Sur \(\mathcal O_{108}\times\mathbb N^2\), le relèvement
\[
 \widetilde F(x,m)=
 \bigl(F_{\rm obs}x,m+c(x\to F_{\rm obs}x)\bigr)
\]
satisfait
\[
 \widetilde{\mathcal K}_{27}(x,m)
 =\bigl(\mathcal K_{27}x,m+(1,18)\bigr)
\]
et
\[
 \widetilde{\mathcal K}_{27}^{\,4}(x,m)
 =\bigl(x,m+(4,72)\bigr)\neq(x,m).
\]
\end{theorem}

\begin{proof}
La constance de \(c_{27}\) sur l'orbite et la loi de concaténation donnent
\[
 c_{108}=\sum_{j=0}^{3}c_{27}(\mathcal K_{27}^jx)
 =4(1,18).
\]
La première coordonnée revient ; la seconde conserve la mémoire accumulée.
\end{proof}

Telle est la donnée qu'une réalisation géométrique doit préserver. On n'identifie pas
la mémoire discrète à la torsion : on construit une représentation
qui permette de les comparer.

\section{Du groupoïde des chemins au groupe de Cartan}

Soit \(\mathcal G_{\rm mem}\) le groupoïde des voies enrichies, avec
composition par concaténation et inverse par renversement. Fixons un élément
\(R_4\in\operatorname{Spin}^+(1,3)\) d'ordre quatre et un vecteur temporel
unitaire \(u\in\mathbb R^{1,3}\) fixé par \(R_4\). Si
\[
c(\gamma)=\bigl(c_g(\gamma),c_s(\gamma)\bigr)\in\mathbb Z^2
\]
est le cocycle additif d'inversion et de changement de feuillet, nous définissons
\[
\rho_{\rm C}^{\rm disc}(\gamma)
=
\left(
R_4^{\,c_g(\gamma)},
\ell_0c_s(\gamma)u
\right)
\in\operatorname{Spin}^+(1,3)\ltimes\mathbb R^{1,3}.
\]

\begin{theorem}[Réalisation discrète du cocycle de retour]
\label{thm:realizacion-discreta-cartan-rev6}
\(\rho_{\rm C}^{\rm disc}\) est un foncteur du groupoïde des voies dans le groupe de
Cartan. En particulier, il conserve l'identité, la concaténation et l'inversion, et
transporte le retour \(27\to54\to108\) sans effacer la coordonnée cumulative.
\end{theorem}

\begin{proof}
La loi du produit semi-direct est
\[
(R,a)(S,b)=(RS,a+Rb).
\]
Comme \(R_4u=u\) et que \(c\) est additif,
\[
\begin{aligned}
\rho_{\rm C}^{\rm disc}(\gamma_2)
\rho_{\rm C}^{\rm disc}(\gamma_1)
&=
\left(
R_4^{c_g(\gamma_2)+c_g(\gamma_1)},
\ell_0\bigl(c_s(\gamma_2)+c_s(\gamma_1)\bigr)u
\right)\\
&=\rho_{\rm C}^{\rm disc}(\gamma_2\gamma_1).
\end{aligned}
\]
La voie identité a un cocycle nul et le renversement change son signe, de sorte
que l'identité et l'inverse sont préservés. Les formules de retour découlent de
\(c_{27}=(1,18)\) et de \(c_{108}=4c_{27}\).
\end{proof}

Ce théorème construit le pont discret. Une réalisation lisse ultérieure
consiste en une représentation
\[
 \rho_{\rm C}:\mathcal G_{\rm mem}
 \longrightarrow\operatorname{Spin}(1,3)\ltimes\mathbb R^{1,3}
\]
et en des champs \(e^I,\omega^{IJ}\) dont l'holonomie entrelace cette représentation.
Le diagramme qui fixe le type est
\[
 \operatorname{Hol}_{\mathcal A_{\ell_0}}
   \bigl(\mathfrak R_{\rm C}(\gamma)\bigr)
 =
 \rho_{\rm C}\bigl(\operatorname{Hol}_{\rm TPK}(\gamma)\bigr).
\]
Les identités, l'inversion, la concaténation et la subdivision doivent être conservées. Cette
équation ne rebaptise pas la retenue courbure : elle exige que le transport
discret déjà représenté et le transport géométrique réalisent la même composition.

Une échelle positive \(\ell_0\) étant fixée, la connexion conjointe est
\[
 \mathcal A_{\ell_0}
 =
 \begin{pmatrix}
 \omega&\ell_0^{-1}e\\
 0&0
 \end{pmatrix}.
\]

\begin{proposition}[Courbure conjointe de Cartan]
\label{prop:curvatura-conjunta-cartan-rev6}
La courbure de \(\mathcal A_{\ell_0}\) est
\[
 \mathcal F_{\ell_0}
 =\dd\mathcal A_{\ell_0}
  +\mathcal A_{\ell_0}\wedge\mathcal A_{\ell_0}
 =
 \begin{pmatrix}
 F_\omega&\ell_0^{-1}T_{\rm tor}\\
 0&0
 \end{pmatrix},
\]
où
\[
 F_\omega^{IJ}
 =\dd\omega^{IJ}+\omega^I{}_K\wedge\omega^{KJ},
 \qquad
 T_{\rm tor}^{I}=D_\omega e^I.
\]
\end{proposition}

\begin{proof}
On multiplie les matrices de formes en utilisant le produit extérieur. Le
bloc diagonal produit \(F_\omega\) ; le bloc de translation produit
\(\dd e+\omega\wedge e=D_\omega e\).
\end{proof}

La courbure et la torsion sont deux composantes d'un même défaut de transport,
mais remplissent des fonctions distinctes. La courbure mesure le retour des repères ;
la torsion mesure la fermeture de la cotétrade.

\section{Gauss--Stokes et capacité surfacique}

Dans un complexe cellulaire fini, pour toute \(1\)-cochaîne \(\Omega\) et toute
\(2\)-chaîne \(z\),
\[
 \boxed{\quad
 \langle\delta\Omega,z\rangle
 =\langle\Omega,\partial z\rangle.
 \quad}
\]
L'identité est la définition du cobord comme adjoint du bord.
Lorsque \(z\) somme les faces orientées de manière cohérente d'une
pseudovariété, les arêtes intérieures s'annulent et seul subsiste le
bord géométrique. La même égalité a déjà été lue comme
vorticité--circulation et comme courbure--holonomie.

L'échelle surfacique possède également un modèle exact. Dans le graphe cubique
\(N\times N\times N\), toute coupe entre deux faces opposées contient au moins
\(N^2\) arêtes et une couche transversale atteint la borne. Pour \(N=9\),
\[
 9^3=729=27^2.
\]
Cette égalité cardinale relie volume et coupe minimale ; elle n'identifie pas les
groupes \((\mathbb Z/9)^3\) et \((\mathbb Z/27)^2\), qui ont des exposants
distincts.

\section{L'action Palatini--Cartan--Holst}

Soit \(M\) une variété lisse orientée de dimension quatre, munie d'une cotétrade non
dégénérée \(e^I\), d'une connexion de Lorentz
\(\omega^{IJ}=-\omega^{JI}\), et de conditions au bord annulant le terme
variationnel de bord. Écrivons
\[
 \Sigma^{IJ}=e^I\wedge e^J,\qquad
 (\star X)^{IJ}=\frac12\epsilon^{IJ}{}_{KL}X^{KL},
 \qquad
 G_{\rm int}=G_{\rm int}(\theta_G),\qquad
 \kappa_{\rm int}=\frac{8\pi G_{\rm int}}{c_{\rm int}^{4}}.
\]
Ces deux dernières quantités ne sont pas définies au moyen de l'action : elles proviennent de
la construction constitutive de la \cref{eq:G-interno}. Le
\cref{thm:homogeneidad-pch} démontre auparavant que la combinaison suivante
appartient à la droite d'action.
L'action du premier ordre est
\[
\boxed{
 S[e,\omega,\Psi]
 =\frac1{2\kappa_{\rm int}c_{\rm int}}\int_M
 \Sigma_{IJ}\wedge
 \left(\star F^{IJ}+\frac1\gamma F^{IJ}\right)
 -\frac{\Lambda_{\rm int}}{\kappa_{\rm int}c_{\rm int}}
  \int_M\operatorname{vol}_e
 +S_{\rm m}[e,\omega,\Psi].}
\]
Le paramètre interne construit au chapitre précédent spécialise
\[
 \gamma=\Gamma_{6\to8}.
\]
La substitution est effectuée après avoir obtenu
\(\Gamma_{6\to8}\) à partir de \(A,C^\ast,q_\pm\) et de l'inclusion
hexade--octade ; l'action n'est pas utilisée pour redéfinir ces entrées.

\begin{theorem}[Équation de Cartan--Holst]
\label{thm:cartan-holst-rev6}
Si le courant de spin est normalisé au moyen de
\[
 \delta_\omega S_{\rm m}
 =-\frac12\int_M\delta\omega^{IJ}\wedge\sigma_{IJ},
\]
la variation par rapport à \(\omega\) produit
\[
 \boxed{\quad
 D_\omega
 \left[(\star+\gamma^{-1}I)\Sigma^{IJ}\right]
 =\kappa_{\rm int}c_{\rm int}\,\sigma^{IJ}.
 \quad}
\]
\end{theorem}

\begin{proof}
On utilise \(\delta F^{IJ}=D_\omega\delta\omega^{IJ}\), on intègre par parties et
on annule le coefficient de la variation arbitraire. Le terme
cosmologique ne dépend pas de \(\omega\).
\end{proof}

En signature lorentzienne, \(\star^2=-I\) sur les bivecteurs. Pour
\(\gamma\in\mathbb R\setminus\{0\}\),
\[
 (\star+\gamma^{-1}I)^{-1}
 =\frac{\gamma^2}{1+\gamma^2}
  (-\star+\gamma^{-1}I).
\]
Par conséquent, si le courant de spin s'annule,
\[
 T_{\rm tor}^I=0,\qquad
 \omega=\mathring\omega(e).
\]
En présence de matière spinorielle, la torsion est déterminée algébriquement par le
courant. Elle ne constitue pas, dans l'action minimale, un degré propagatif
indépendant.

\section{Équation métrique, spin et bilan géométrique}

Pour fixer sans ambiguïté la normalisation, nous définissons les sources mesurées dans
la droite d'action par
\[
 \delta_g S_{\rm m}
 =-\frac12\int_M\operatorname{vol}_e\,
 \mathcal T^{S}_{\mu\nu}\,\delta g^{\mu\nu}.
\]
Lorsque les coordonnées de la tétrade ont la dimension d'une longueur, nous écrivons
\[
 T^{\rm phys}_{\mu\nu}:=c_{\rm int}\mathcal T^{S}_{\mu\nu},
 \qquad
 U^{\rm phys,spin}_{\mu\nu}
 :=c_{\rm int}\mathcal U^{S,\rm spin}_{\mu\nu}.
\]
Après élimination algébrique de la connexion, l'équation métrique adopte les
deux formes équivalentes
\[
 \boxed{
 G_{\mu\nu}+\Lambda_{\rm int}g_{\mu\nu}
 =\kappa_{\rm int}c_{\rm int}
 \bigl(\mathcal T_{\mu\nu}^{S,\rm LC}
       +\mathcal U_{\mu\nu}^{S,\rm spin}\bigr)
 =\kappa_{\rm int}
 \bigl(T_{\mu\nu}^{\rm phys,LC}
       +U_{\mu\nu}^{\rm phys,spin}\bigr).}
\]
La contribution de spin est quadratique en le courant correspondant et
dépend du couplage de matière ; son coefficient n'est pas fixé sans déclarer
l'action fermionique et les conventions de dualité.

Les identités géométriques
\[
 D_\omega T_{\rm tor}^{I}=F^{I}{}_{J}\wedge e^{J},
 \qquad
 D_\omega F^{IJ}=0
\]
impliquent, après élimination de la torsion,
\[
 \nabla^\mu\bigl(T_{\mu\nu}^{\rm phys,LC}
             +U_{\mu\nu}^{\rm phys,spin}\bigr)=0.
\]
Ce bilan de Bianchi--Noether est une condition de cohérence du
transducteur gravitationnel : aucun terme supplémentaire ne peut être incorporé sans
provenir d'une action compatible ou satisfaire le même bilan.

\section{Correspondance surfacique–volumétrique \(\pi/\varphi^2\)}

Le calendrier \((\Sigma,\Pi,\Pi)\) produit la matrice d'incidence
\[
 F_{\rm av}=\begin{pmatrix}0&1\\1&1\end{pmatrix}.
\]
Sa mesure de Parry donne
\(\mu_{\rm ar}/\mu_{\rm vol}=\varphi^{-2}\), et le lecteur régional marque une
seule fois la fermeture \(\pi_{\rm HMT}\). Le retour marqué satisfait
\[
 \Theta_n
 =\pi_{\rm HMT}\frac{F_{n-1}}{F_{n+1}}
 \longrightarrow
 \frac{\pi_{\rm HMT}}{\varphi^2}.
\]
Ce rapport n'est pas une autre version de \(G\). C'est la correspondance adimensionnelle entre
lecture surfacique et lecture volumétrique. Sous une même amplitude positive
\(\mathcal M\),
\[
 \mathcal I_{\rm av}=\mu_{\rm vol}\mathcal M,\qquad
 \mathcal G_{\rm av}=\pi_{\rm HMT}\mu_{\rm ar}\mathcal M
\]
donnent exactement
\[
 \frac{\mathcal G_{\rm av}}{\mathcal I_{\rm av}}
 =\frac{\pi_{\rm HMT}}{\varphi^2}.
\]
Sur le feuillet plan, la restriction diagonale est
\(\mathcal G_{\rm tor}=\mathcal I_{\rm tor}\). L'équivalence locale et
l'ambivalence de bord sont donc deux régimes du lecteur, non deux
valeurs incompatibles d'une masse.

\section{Typage de \(G\) comme transducteur gravitationnel et ordre causal de ses lecteurs}

Le type et l'homogénéité de \(G\) appartiennent à la construction
constitutive préalable. La \cref{eq:G-interno} fixe une seule fois la famille
covariante \(G_{\rm int}(\theta_G)\), et le
\cref{thm:homogeneidad-pch} démontre que son insertion dans l'action possède
la dimension d'une action. Ce chapitre ne redéfinit pas cette famille : il utilise sa
section avec l'holonomie déjà construite. La causalité n'est pas linéaire,
mais convergente :
\[
\begin{aligned}
 \mathcal B_{\rm geom}&:
 \text{retour et cocycle}\to\text{représentation de Cartan}
 \to(e,\omega),\\
 \mathcal B_{\rm dim}&:
 (\hbar_{\rm int},c_{\rm int},t_0,\theta_G)
 \to G_{\rm int}\to\kappa_{\rm int},
\end{aligned}
\]
\[
 (\mathcal B_{\rm geom},\mathcal B_{\rm dim})
 \to S_{\rm PCH}
 \to\text{torsion algébrique, équation métrique et bilan}.
\]
La comparaison métrologique intervient à la fin et ne sélectionne aucune des
deux branches.

\Needspace{22\baselineskip}
\section{Origine collective du mode gravitationnel dans l'holonomie des voies}

L'action précédente a pour objets fondamentaux la cotétrade et la connexion.
La torsion de Cartan est algébrique ; la courbure possède le secteur propagatif qui
contient les ondes gravitationnelles. Le formalisme n'a donc pas besoin de
postuler un graviton élémentaire pour définir l'interaction. Une éventuelle
quantification des perturbations peut produire un mode collectif
d'hélicités \(\pm2\) ; ce mode serait une excitation du transport
géométrique, non une donnée ajoutée au niveau d'APP, de TRIT ou de TPK.

Cette distinction exclut également une confusion fréquente. L'existence
d'ondes classiques démontre une propagation de courbure ; elle ne démontre pas à elle seule
un quantum fondamental. Réciproquement, décrire le champ au moyen de la
géométrie n'élimine pas ses degrés radiatifs. HMT--MD situe la question au
bon endroit : on construit d'abord la mémoire et son holonomie ; on étudie ensuite
le spectre de ses excitations.

\end{HMTScientificOwnerSubordinated}

\paragraph{Comparaison typée du quantum d’aire avec l’énergie, l’entropie et l’information d’horizon.} L’action du premier ordre permet de comparer le quantum aréal à l’énergie, à l’entropie et à l’information d’horizon sans confondre ces observables.
